\documentclass[a4paper,12pt,twoside, titlepage]{book}
\usepackage[latin,french]{babel}

\usepackage{fontspec}
\defaultfontfeatures[Old Standard]{%utiliser oldstandard pour imprimer le grec
Extension = .otf, 
UprightFont = OldStandard-Regular,  
ItalicFont = OldStandard-Italic,  
BoldFont = OldStandard-Bold,  
BoldItalicFont = OldStandard-BoldItalic,  
SlantedFont = OldStandard-Regular,
SlantedFeatures = {FakeSlant=0.25},
BoldSlantedFont = OldStandard-Bold,
BoldSlantedFeatures = {FakeSlant=0.25}
}
\setmainfont{Old Standard}
\usepackage[poetry=verse]{ekdosis}

\SetAlignment{
texts=translation;edition,
apparatus=edition
}

\title{exercice Ekdosis}
\author{Clara Grometto}
\date{19/04/2026}

%Babel gère les règles de césure et de ponctuation différemment selon la langue. Ces deux lignes font commuter automatiquement la langue à chaque bloc, évitant de le faire manuellement dans chaque edition.

\AtBeginEnvironment{edition}{\selectlanguage{latin}}
\AtBeginEnvironment{translation}{\selectlanguage{french}}

%pour cet exercice, s'appoyer sur la photographie du CS, sur les facsimilés des deux pages à encoder, ainsi que sur le modèle pdf simplifié

%1. Déclarer les sources avec \DeclareWitness
% syntaxe : \DeclareWitness{identifiant}{sigle affiché}{nom complet}
%           [clé=valeur] pour les options
%L'identifiant (1er argument) est utilisé dans le code LaTeX pour référencer ce manuscrit. Le sigle affiché (2e arg.) est ce qui apparaît imprimé. Le nom complet sert à la liste des témoins.

%compléter : finir d'encoder le Conspectus Siglorum. On omet ici l'encodage des mains successives. Mais on pourrait les mettre dans un \DeclareHand{M1}{M}{M\textsuperscript{1}}[Codicis corrector]

\DeclareWitness{M}{\emph{M}}{Codex Mediceus}
[origDate={saec. \textsc{v}}]
\DeclareWitness{}{}{}
  [origDate={}]
\DeclareWitness{}{}{}
  [origDate={}]
\DeclareWitness{}{}{}
  [origDate={}]

%2. déclarer un shorthand pour les codex P et R
%syntaxe : ........
%ce sont des sources qui sont souvent citées ensemble car elles présentent souvent la même leçon

\DeclareShorthand{}{}{}

\DeclareWitness{a}{\emph{α}}{Codex Mediolanensis biblioth. Ambrosianae, olim Petrarcae} %les id en alphabet latin car c'est trop compliqué de les passer en grec
  [origDate={saec. \textsc{vxiii-xiv}}]
\DeclareWitness{c}{\emph{γ}}{Codex Guelferbytanus Gudianus 2\textsuperscript{o}, 70.}
  [origDate={saec. \textsc{ix}}]


%3. Déclarer les sources
%syntaxe : \DeclareSource{identifiant}{sigle affiché}
%Les sources sont des textes anciens (commentaires, citations) qui attestent une leçon du texte mais ne sont pas des manuscrits du texte lui-même.
%compléter : finir de déclarer les sources

\DeclareSource{Serv}{\textsc{Serv.}} %@MV: Servius (4e s. apr.) a commenté les Bucoliques de Virgile (https://www.perseus.tufts.edu/hopper/text?doc=Perseus%3atext%3a2007.01.0092). Donc pour établir le texte des Buc. de Virg., on consulte aussi ce genre de textes anciens. C'est une source. 
\DeclareSource{}{} 
\DeclareSource{}{}
\DeclareSource{}{} 
\DeclareSource{}{}

\DeclareSource{Quintil}{\textsc{Quintil.}} %Quintil et Asper ne sont pas recensés dans le CS mais ils pourraient être considérés comme des sources. 
\DeclareSource{Asper}{\textsc{Asper.}}

\DeclareSource{havet}{\emph{Havet}} %havet est une source contemporaine. On aurait alors un fichier .bib avec toutes les informations bibliographique, et havet serait la clé

\SetLineation{lineation=none,vmodulo} %numérotation toutes les 5 lignes, plus difficile à implémenter pour la version française. 

\DeclareApparatus{default}

\begin{document}


%afficher la page de titre


%ekdosis organise le texte en blocs synchronisés. Voici comment les comprendre.
%Voici le schéma à suivre pour les environnements imbriqués
%\begin{alignment}        zone synchronisée

%\begin{translation}     traduction française
% ...texte français...
%\end{translation}

%\begin{edition}         texte latin
%...texte latin...
%\end{edition}

%\end{alignment}

%Chaque alignment correspond à un paragraphe : traduction et édition sont alignées ensemble. Si on change de locuteur, on ouvre un nouveau bloc alignment.


%Voici les titre : 
%- ouvrir un environnement alignment 
%- dans cet environnement placer deux autres environnement translation et edition

\begin{center}\huge\textsc{Virgile}\end{center}

\begin{center}\large\textsc{Bucoliques}\end{center}

\begin{center}\large I\end{center}



\begin{center}\huge\textsc{P. Vergili Maronis}\end{center}

\begin{center}\large\textsc{Bvcolica}\end{center}

\begin{center}\large I\end{center}

%voici les textes: 
%- ouvrir un nouvel environnement alignment 
%- dans cet environnement, ouvrir autant de paires translation et edition que de strophes. 
%- les vers sont enveloppés d'un environnement ekdverse
%- encoder l'apparat critique

%%% --- PAR. 1 --- %%%

\begin{center}\textsc{Mélibée}\end{center}

Toi, Tityre, étendu sous le couvert d'un large hêtre, tu essaies un air silvestre sur un mince pipeau; nous autres, nou quittons notre pays et nos chères campagnes; loin du pays nous sommes exilés; toi, Tityre, nonchalant sous l'ombrage, tu apprends aux bois à redire le nom de la belle Amaryllis. 

\begin{center}\textsc{Meliboevs}\end{center}

Tityre, tu patulae recubans sub tegmine fagi\\
\app{}
tenui musa meditaris auena~;\\
nos patriae finis et dulcia linquimus arua~;\\
nos patriam fugimus~; tu, Tityre, lentus in umbra,\\
\app{} resonare doces Amaryllida siluas.\\

%%% --- PAR. 2 --- %%%

\begin{center}\textsc{Tityre}\end{center}

O Mélibée, c'est à un dieu que nous devons ces loisirs~; car Il sera pour moi toujours un dieu~; son autel, une tendre victime, un agneau de nos bergeries, souvent l'ensanglantera. Grâce à Lui, mes génisses ont le droit de paître en liberté, comme tu vois, et moi-même celui de jouer mes airs préférés sur un roseau rustique.

\begin{center}\textsc{Tityrvs}\end{center}

O Meliboee, deus nobis haec otia fecit~:\\
namque erit ille mihi semper deus~; illius aram\\
saepe tener nostris ab ouilibus imbuet agnus.\\
Ille meas errare boues, ut cernis, et ipsum\\
ludere quae uellem calamo permisit agresti.\\

%%% --- PAR. 3 --- %%%

\begin{center}\textsc{Mélibée}\end{center}

Je ne suis point jaloux, mais étonné plutôt~; partout, dans les campagnes, il y a tant de désordre~! Vois~: mes pauvres chèvres, je les pousse, dolent, droit devant moi~; et celle-ci, Tityre, j'ai peine à la tirer~: ici, dans une touffe de coudriers, elle vient de mettre bas deux bessons, espoir du troupeau, hélas~! laissés sur le roc nu.

\begin{center}\textsc{Meliboevs}\end{center}

Non equidem inuideo, miror magis~: undique totis\\
usque adeo \app{
} agris~! En ipse capellas\\
\app{
} aeger ago~; hanc etiam uix, Tityre, duco~:\\
hic inter densas corylos modo namque gemellos,\\
spem gregis, a~! silice in nuda conixa reliquit.\\

\end{document}



